% Adición junto a eq:alpha-recurrencia o inc-ampl-normalizacion.
% No cambia ni elimina ninguna fórmula ni renombra sus símbolos.
\paragraph{Balance firmado y suma con cociente entrante.}
\label{ampl-alfa-z-s-aclaracion}
Conviene distinguir las dos etapas consecutivas de la división
posicional mediante una notación común. El balance anterior al transporte
del cociente es
\[
 Z_j:=P_j+E_j-\Phi_j-K_j^{\mathrm{per}}.
\]
La variable $s_j$ de \eqref{eq:alpha-recurrencia} designa la suma que ya
incluye la contribución entrante, por lo que
\[
 s_j=Z_j+c_{j+1},\qquad
 A_j=Z_j+c_{j+1}-1000c_j.
\]
En \eqref{inc-ampl-normalizacion}, la letra $s_j$ designa, localmente,
el balance anterior a esa contribución, es decir, $Z_j$ en la notación
presente. Las dos fórmulas expresan, por tanto, la misma división
euclídea con etapas nominales diferentes. La configuración incidencial
$H^-_\alpha$ utiliza los signos del balance $Z_j$ de la ventana inicial;
la determinación de las cifras utiliza además los cocientes de enlace
$c_{j+1}$ y $c_j$. Esta distinción conserva conjuntamente el soporte
firmado y la normalización de la expansión.
